% GENERATED FILE. Edit canonical lessons, then run npm run generate:books.
% canonical-source-hash: fnv1a64:5763eee6373c3812
% canonical-lessons: IT-C459-prelevare, IT-C459-prelievo, IT-C459-addebitare, IT-C459-bonifico, IT-C459-rata

\chapter{In banca}
\label{ch:it-in-banca}
\hlchaptermodality{459}

\begin{chapteropening}
\textbf{By the end of this chapter:} \emph{I can manage money at the bank: withdrawing cash, a withdrawal, being charged, a bank transfer and instalments.}

Complete the last lesson of chapter 459: I can manage money at the bank: withdrawing cash, a withdrawal, being charged, a bank transfer and instalments.
\end{chapteropening}

\section[prelevare]{prelevare — to withdraw (money)}

\label{lesson:IT-C459-prelevare}



\begin{quote}
Before the new one: say the Italian for to put postage on, then the Italian for to deliver.
\end{quote}

\subsection*{You'll want to know: prelevare}
\textbf{prelevare} — \textquotedblleft{}to withdraw (money)\textquotedblright{}.

\begin{quote}
Devo prelevare al bancomat.
\end{quote}

\begin{quote}
\emph{I need to withdraw cash at the cash machine.}
\end{quote}

\emph{Prelevare il sangue} — to take a blood sample.

\begin{cousinweb}
\emph{prelevare} is a learned borrowing from Late Latin \emph{praelevāre} 'to lift first', from \emph{levāre} 'to lift' (Lewis \& Short).
\end{cousinweb}

\begin{grammarlens}[title={What you've built}]
One more word for telling, explaining and describing at B1.
\end{grammarlens}

\subsection*{Your turn}
\begin{itemize}
\raggedright
  \item \emph{Say it:} \emph{prelevare}
  \item \emph{Say it:} \emph{prelevare}, once more
  \item \emph{Say it:} a sentence of your own with \emph{prelevare}
  \item \emph{Recall:} say the Italian for to put postage on, then the Italian for to deliver, then say \emph{prelevare} again
\end{itemize}

\begin{tcolorbox}[breakable,colback=teal!4,colframe=teal!35!black,title={Before you move on}]
What does prelevare mean? (\textquotedblleft{}To withdraw (money)\textquotedblright{}.) Say it once more.
\end{tcolorbox}
\section[il prelievo]{il prelievo — the withdrawal; the sample (of blood)}

\label{lesson:IT-C459-prelievo}



\begin{quote}
Before the new one: say the Italian for to deliver, then the Italian for to withdraw.
\end{quote}

\subsection*{You'll want to know: il prelievo}
\textbf{il prelievo} — \textquotedblleft{}the withdrawal; the sample (of blood)\textquotedblright{}.

\begin{quote}
Il prelievo massimo è di cinquecento euro.
\end{quote}

\begin{quote}
\emph{The maximum withdrawal is five hundred euros.}
\end{quote}

At the doctor's, \emph{il prelievo} is a blood test.

\begin{cousinweb}
\emph{prelievo} is built on \emph{prelevare} 'to withdraw', a learned borrowing from Late Latin \emph{praelevāre} 'to lift first', from \emph{levāre} 'to lift'.
\end{cousinweb}

\begin{grammarlens}[title={What you've built}]
One more word for telling, explaining and describing at B1.
\end{grammarlens}

\subsection*{Your turn}
\begin{itemize}
\raggedright
  \item \emph{Say it:} \emph{il prelievo}
  \item \emph{Say it:} \emph{il prelievo}, once more
  \item \emph{Say it:} a sentence of your own with \emph{il prelievo}
  \item \emph{Recall:} say the Italian for to deliver, then the Italian for to withdraw, then say \emph{il prelievo} again
\end{itemize}

\begin{tcolorbox}[breakable,colback=teal!4,colframe=teal!35!black,title={Before you move on}]
What does il prelievo mean? (\textquotedblleft{}The withdrawal; the sample (of blood)\textquotedblright{}.) Say it once more.
\end{tcolorbox}
\section[addebitare]{addebitare — to charge (to an account), to debit}

\label{lesson:IT-C459-addebitare}



\begin{quote}
Before the new one: say the Italian for to withdraw, then the Italian for the withdrawal.
\end{quote}

\subsection*{You'll want to know: addebitare}
\textbf{addebitare} — \textquotedblleft{}to charge (to an account), to debit\textquotedblright{}.

\begin{quote}
La banca mi ha addebitato due euro.
\end{quote}

\begin{quote}
\emph{The bank charged me two euros.}
\end{quote}

\emph{L'addebito} — the debit.

\begin{cousinweb}
\emph{addebitare} is built from \emph{a-} and \emph{debito} 'debt', a learned borrowing from Latin \emph{dēbitum} 'what is owed', from \emph{dēbēre} 'to owe' (Lewis \& Short).
\end{cousinweb}

\begin{grammarlens}[title={What you've built}]
One more word for telling, explaining and describing at B1.
\end{grammarlens}

\subsection*{Your turn}
\begin{itemize}
\raggedright
  \item \emph{Say it:} \emph{addebitare}
  \item \emph{Say it:} \emph{addebitare}, once more
  \item \emph{Say it:} a sentence of your own with \emph{addebitare}
  \item \emph{Recall:} say the Italian for to withdraw, then the Italian for the withdrawal, then say \emph{addebitare} again
\end{itemize}

\begin{tcolorbox}[breakable,colback=teal!4,colframe=teal!35!black,title={Before you move on}]
What does addebitare mean? (\textquotedblleft{}To charge (to an account), to debit\textquotedblright{}.) Say it once more.
\end{tcolorbox}
\section[il bonifico]{il bonifico — the bank transfer}

\label{lesson:IT-C459-bonifico}



\begin{quote}
Before the new one: say the Italian for the withdrawal, then the Italian for to charge.
\end{quote}

\subsection*{You'll want to know: il bonifico}
\textbf{il bonifico} — \textquotedblleft{}the bank transfer\textquotedblright{}.

\begin{quote}
Pago la caparra con un bonifico.
\end{quote}

\begin{quote}
\emph{I pay the deposit by bank transfer.}
\end{quote}

\emph{Fare un bonifico} — to make a transfer.

\begin{cousinweb}
\emph{bonifico} is built on \emph{bonificare} 'to credit, to improve', a learned formation from Latin \emph{bonus} 'good' and \emph{facere} 'to make'.
\end{cousinweb}

\begin{grammarlens}[title={What you've built}]
One more word for telling, explaining and describing at B1.
\end{grammarlens}

\subsection*{Your turn}
\begin{itemize}
\raggedright
  \item \emph{Say it:} \emph{il bonifico}
  \item \emph{Say it:} \emph{il bonifico}, once more
  \item \emph{Say it:} a sentence of your own with \emph{il bonifico}
  \item \emph{Recall:} say the Italian for the withdrawal, then the Italian for to charge, then say \emph{il bonifico} again
\end{itemize}

\begin{tcolorbox}[breakable,colback=teal!4,colframe=teal!35!black,title={Before you move on}]
What does il bonifico mean? (\textquotedblleft{}The bank transfer\textquotedblright{}.) Say it once more.
\end{tcolorbox}
\section[la rata]{la rata — the instalment}

\label{lesson:IT-C459-rata}



\begin{quote}
Before the new one: say the Italian for to charge, then the Italian for the bank transfer.
\end{quote}

\subsection*{You'll want to know: la rata}
\textbf{la rata} — \textquotedblleft{}the instalment\textquotedblright{}.

\begin{quote}
Pago la macchina a rate.
\end{quote}

\begin{quote}
\emph{I am paying for the car in instalments.}
\end{quote}

\emph{A rate} — in instalments.

\begin{cousinweb}
\emph{rata} is inherited from Latin \emph{rata} (EtymDB), short for \emph{pars rata} 'a fixed share', from \emph{rērī} 'to reckon'.
\end{cousinweb}

\begin{grammarlens}[title={What you've built}]
One more word for telling, explaining and describing at B1.
\end{grammarlens}

\subsection*{Your turn}
\begin{itemize}
\raggedright
  \item \emph{Say it:} \emph{la rata}
  \item \emph{Say it:} \emph{la rata}, once more
  \item \emph{Say it:} a sentence of your own with \emph{la rata}
  \item \emph{Recall:} say the Italian for to charge, then the Italian for the bank transfer, then say \emph{la rata} again
\end{itemize}

\begin{tcolorbox}[breakable,colback=teal!4,colframe=teal!35!black,title={Before you move on}]
What does la rata mean? (\textquotedblleft{}The instalment\textquotedblright{}.) Say it once more.
\end{tcolorbox}
